
\label{sec:dexterity}

\noindent \textbf{\MethodName{} achieves high performance on dexterous tasks without task-specific post-training.}
In our first set of experiments, we study how well \MethodName{} can master dexterous tasks that were seen in the training data, but where the goal is to perform these tasks as robustly and efficiently as possible. This is surprisingly difficult for prior robotic foundation models: often the best-performing policies are fine-tuned for specific downstream tasks, even if they use generalist pre-training~\citep{pistar06, pi06model}. We aim to answer: can the general-purpose \MethodName{} model match the performance of task-specific fine-tuned models on a variety of dexterous manipulation tasks?

We use the tasks shown in Fig.~\ref{fig:distillation_results}. These include the espresso making, box building, and laundry folding tasks that we previously used to evaluate the RL-trained \Piszs{} models~\citep{pistar06}, where we can directly compare the speed and robustness of the single general-purpose \MethodName{} model to the individual RL-finetuned specialist \Piszs{} models. We also study a number of other dexterous tasks, including some tasks from our previous ``Robot Olympics'' experiments (making a peanut butter sandwich, turning a shirt inside-out, and driving through a door) and a number of additional dexterous tasks, such as fully slicing up a zucchini, peeling a few fruits and vegetables (zucchini, cucumbers, and carrots), and a long-horizon task that involves replacing a trash bag in a trash can. We find that \MethodName{} achieves performance that is competitive with the RL specialists used in the \Piszs{} release \citep{pistar06} for all of the tasks considered in the paper directly out of the box (Fig.~\ref{fig:distillation_results}, first row), and even outperform the specialists in throughput in the difficult laundry and box building tasks. Additionally, we compare \MethodName{} to SFT specialists trained on top of \Pizs{} for a number of other dexterous tasks, and find that \MethodName{} is again able to closely match the performance of all specialist policies (Fig.~\ref{fig:distillation_results}, second row).

To understand how the training recipe of \MethodName{} affects performance, we additionally compare \MethodName{} with two ablations on the tasks from the \Piszs{} release: \MethodName{} (no eval data), which holds out all autonomous evaluation episodes from training (and thus cannot benefit from distilling agent rollouts from strong, e.g. RL trained, policies), and \MethodName{} (no metadata), which omits episode metadata from the context as shown in Fig.~\ref{fig:distillation_ablations} for the tasks used in the \Piszs~ release. Results suggest that \MethodName{} significantly outperforms both \MethodName{} (no eval data) and \MethodName{} (no metadata) on all tasks. Since policy evaluation data can vary widely in quality, training on this data, combined with rich metadata to disambiguate high and low quality behaviors, is critical for \MethodName{}'s strong performance on all of these challenging tasks.

\noindent \textbf{\MethodName{} achieves high performance on tasks that require memory without fine-tuning.}
In these experiments, we study how well \MethodName{} can perform tasks that require explicitly keeping track of previous observations~\citep{torne2026mem}. 
We compare the same single \MethodName{} model out of the box to the task-specific fine-tuned versions of \Pizs{} with memory used in~\citet{torne2026mem}, and find that \MethodName{} can achieve similar or better performance to the fine-tuned specialists on all of these tasks (Fig.~\ref{fig:memory}).

